% Pursuing Stacks (cotes 134-2 to 134-8): archive cote 134-7, site page 8 (PDF page 9), typescript page 450.
% typescript. Draft compared: ps6.tex:329-399.
% AI-transcribed by Claude Opus 5.5 (claude-opus-5-5, adaptive thinking, effort high), Message Batches API, batch ps3 (msgbatch_01XdFmuVVBAcgc7EFxdtCqsj), 2026-11-12; not yet checked by eye.
\scanpage{134-7/8}
\mspage{450}

\diagram{\begin{tikzcd}
Hot_n(k) \arrow[r] \arrow[d] & D.(Ab_k) \arrow[d] \\
Hot_n \arrow[r] & D\hand{.}(Ab)
\end{tikzcd}}
\margin{(l)}
,

where the second horizontal arrow is the usual abelianization functor
for homotopy types,  and the second vertical one comes from ring restric-%
tion functor  $Ab_k \longrightarrow Ab = Ab_{\underline{Z}}$ .  In \hand{(}k'\hand{),} Ch.  denotes the category
of chain complexes,  and it looks rather mysterious again how to get
such a functor\hand{(}k').  \xout{But} We may of course think of the trivial abeliani-%
zation
\[ X. \mapsto k^{(X.)} \stackrel{def}{=} (n \mapsto k^{(X_n)} ) \quad , \]
where for an object X in U(k),  or more generally any sheaf on $Aff_{/k}$ ,
$k^{(X)}$ defines trivial k-linearization of this sheaf,  in the sense of the
topos of all such sheaves.  Anyhow $k^{(X)}$ is a sheaf,  not a k-module,
so we should still take sections to get what we want  -  but this \inserted{functor} looks
not only prohibitively large and unaccessible,  but just silly\hand{!}  A much
better choice for k-linearizing objects of U(k) specifically seems the
following.  Disregarding the fibration structures,  such an object X is
isomorphic to an object W(M),  M some k-module.  We look for a k-lineari-%
zation
\[ \text{(m)} \qquad X \hand{(}\over{=}{\simeq}W(M)) \longrightarrow W(L(\hand{X})) \quad , \]
where L(X) is a suitable k-module.  Now,  among all maps
\[ X \longrightarrow W(N) \]
\margin{X} of X into sheaves \uline{of the type} W(N),  there is a universal one,  which
in terms of M can be described as
\[ N = \hand{\Gamma}_k(M) \quad , \]
where $\hand{\Gamma}_k$  denotes the ``algebra with devided powers\deleted{''} generated by M\hand{''},
the canonical map  $M \longrightarrow \hand{\Gamma}_k(M)$  or rather
\[ \hand{(}\text{n}) \qquad W(M) \longrightarrow W(\hand{\Gamma}_k(M)) \quad , \quad x \mapsto exp(x) = \hand{\sum_{i\geq 0}} x^{(i)} \]
being \uline{the} ``\uline{universal}\deleted{''} polynomial map\hand{''}  of M with values in a module \over{M}{N}
(or rather,  of W(M) into W(N)).   Here  $x^{(i)}$ denotes the i.th devided
power of x ,  which is an element of  $\hand{\Gamma}_k^{\,i}(M)$ .  It just occurs to me
that this expression of exp(x),  the universal map,  is infinite,  thus
it does'nt take its values in W(N)  actually,  but in a suitable comple-%
tion of it - this does'nt seem too serious a drawback though !  The
point I wish to make here,  is that for given X in U(k),  defining

\margin{8}
